\documentclass[12pt,a4paper]{article}
\usepackage{fontspec}
%\setmainfont{Times New Roman}
\usepackage[margin=1in,left=1in,top=1in,includefoot]{geometry}
\usepackage{fancyhdr}
\usepackage{ragged2e}
\usepackage{blindtext}
\usepackage{setspace}
\usepackage{graphicx}
% <-- Activate global fancy border
\begin{document}
\newpage

\pagenumbering{arabic}
\pagestyle{fancy}
\fancyhead[R]{\textit{Automatic Tree Climber}}       % remove subsection name
\fancyhead[L]{}  % your header text
\fancyfoot{}
\fancyfoot[L]{Dept. of CSE, Vemana IT}
\fancyfoot[C]{\thepage}
\fancyfoot[R]{2026-27} 
\renewcommand{\footrulewidth}{1pt}

\begin{flushleft}
{\fontsize{16pt}{16pt}\selectfont \textbf{CHAPTER 1}}
\end{flushleft}		
\begin{center}
{\fontsize{18pt}{18pt}\selectfont \textbf{INTRODUCTION}}\\[1mm]
\end{center}
\onehalfspacing
%\blindtext  \\
\justifying \noindent The project unveils in helping fire accidents which can be taken place anywhere and anytime. Fire accidents have been a persistent threat to human safety, causing extensive damage to life and property. Traditional firefighting techniques involve human firefighters, which can be hazardous, time-consuming, and sometimes ineffective in controlling large-scale fires. To overcome these challenges, technological advancements in robotics and automation have paved the way for innovative solutions in fire detection and suppression. \\[2mm]
\begin{flushleft}
{\fontsize{14}{20}\selectfont\textbf{\*1.1 PROBLEM STATEMENT}}\\[3mm]
\justifying \noindent The Hybrid Fire Fighting Robot is designed to autonomously detect and extinguish fires while 
minimizing human risk. It is equipped with flame, smoke, and thermal sensors for accurate fire detection, along with a high-pressure water pump for effective suppression. \\[2mm]
\end{flushleft}
\begin{flushleft}
{\fontsize{14}{20}\selectfont\textbf{\*1.2 SCOPE}}\\[3mm]
\justifying \noindent The Hybrid Fire Fighting Robot is designed to autonomously detect and extinguish fires while 
minimizing human risk. It is equipped with flame, smoke, and thermal sensors for accurate fire detection, along with a high-pressure water pump for effective suppression. \\[2mm]
\end{flushleft}
%\blindtext
\begin{flushleft}
{\fontsize{14}{20}\selectfont\textbf{\*1.3 LIST OF OBJECTIVES} }	\\[3mm]
\justifying \noindent The primary objective of the Hybrid Fire Fighting Robot is to develop an autonomous and remotely 
controlled fire suppression system that enhances safety and efficiency in fire emergencies. The 
robot is designed to detect fires using flame, smoke, and thermal sensors and respond by navigating 
toward the fire source and extinguishing it with a high-pressure water pump. 	\\[2mm]
\end{flushleft}

\begin{flushleft}
{\fontsize{12}{20}\selectfont\textbf{\*1.3.1 Plan of Action for the project \small{(timeline chart)}} }	\\[3mm]
\justifying \noindent The primary objective of the Hybrid Fire Fighting Robot is to develop an autonomous and remotely 
controlled fire suppression system that enhances safety and efficiency in fire emergencies. The 
robot is designed to detect fires using flame, smoke, and thermal sensors and respond by navigating 
toward the fire source and extinguishing it with a high-pressure water pump. 	

\end{flushleft}		
%\blindtext
\newpage
\begin{flushleft}
{\fontsize{16pt}{18pt}\selectfont \textbf{CHAPTER 2}}
\end{flushleft}		
\begin{center}
{\fontsize{18pt}{22pt}\selectfont \textbf{LITERATURE SURVEY}}
\end{center}
\onehalfspacing
\blindtext\\
\begin{flushleft}
{\fontsize{12}{14}\selectfont\textbf{\cite{1} Title:
\newline Author:
\newline Year:}	}								
\end{flushleft}		
\blindtext\\
{\fontsize{12}{14}\selectfont\textbf{Concerns:
}}								
\begin{flushleft}
{\fontsize{12}{14}\selectfont\textbf{\cite{2}Title:
\newline Author:
\newline Year:}		}								
\end{flushleft}		
\blindtext\\
{\fontsize{12}{14}\selectfont\textbf{Concerns:
}}								

\begin{flushleft}
{\fontsize{12}{14}\selectfont\textbf{\cite{3}Title:
\newline Author:
\newline Year:}		}								
\end{flushleft}		
\blindtext\\
{\fontsize{12}{14}\selectfont\textbf{Concerns:
}}								
\begin{flushleft}
{\fontsize{12}{14}\selectfont\textbf{\cite{4}Title:
\newline Author:
\newline Year:}		}								
\end{flushleft}		
\blindtext\\
{\fontsize{12}{14}\selectfont\textbf{Concerns:
}}								

\begin{flushleft}
{\fontsize{12}{14}\selectfont\textbf{\cite{5}Title:
\newline Author:
\newline Year:}		}								
\end{flushleft}		
\blindtext\\
{\fontsize{12}{14}\selectfont\textbf{Concerns:
}}								
\begin{flushleft}
{\fontsize{12}{14}\selectfont\textbf{\cite{6}Title:
\newline Author:
\newline Year:}		}								
\end{flushleft}		
\blindtext\\
{\fontsize{12}{14}\selectfont\textbf{Concerns:
}}	
\begin{flushleft}
{\fontsize{14}{20}\selectfont\textbf{\*2.1 \MakeUppercase{Comparative Analysis (In table format)} }}	\\[3mm] 	
\end{flushleft}
\newpage
\begin{flushleft}
{\fontsize{16}{20}\selectfont \textbf{CHAPTER 3}}
\end{flushleft}		
\begin{center}
{\fontsize{18}{18}\selectfont \textbf{SYSTEM REQUIREMENTS}}
\end{center}
\onehalfspacing
\begin{flushleft}
{\fontsize{14}{20}\selectfont\textbf{\*3.1 Functional Requirements}		}											
\end{flushleft}		
\blindtext

\begin{flushleft}
{\fontsize{14}{20}\selectfont\textbf{\*3.2 Non Functional Requiements}		}											
\end{flushleft}		
\blindtext
\begin{flushleft}
{\fontsize{14}{20}\selectfont\textbf{\*3.3 Hardware Requirements}		}											
\end{flushleft}		
\blindtext
\\
\begin{flushleft}
{\fontsize{12}{12}\selectfont\textbf{\*3.3.1 SUB SECTION 3}		}											
\end{flushleft}		
\blindtext

\begin{flushleft}
{\fontsize{14}{20}\selectfont\textbf{\*3.4 Software Requirements}		}											
\end{flushleft}		
\blindtext

\begin{flushleft}
{\fontsize{12}{12}\selectfont\textbf{\*3.4.1 SUB SECTION 1}		}											
\end{flushleft}		
\blindtext

\begin{flushleft}
{\fontsize{12}{12}\selectfont\textbf{\*3.4.2 SUB SECTION 2}		}											
\end{flushleft}		
\blindtext
\newpage
\onehalfspacing
\begin{flushleft}
{\fontsize{16pt}{16pt}\selectfont \textbf{CHAPTER 4}}
\end{flushleft}		
\begin{center}
{\fontsize{18pt}{22pt}\selectfont \textbf{DESIGN METHODOLOGY}}
\end{center}
\begin{flushleft}
{\fontsize{14}{20}\selectfont\textbf{\*4.1 System Architecture}		}											
\end{flushleft}		
\blindtext

\begin{flushleft}
{\fontsize{14}{20}\selectfont\textbf{\*4.2 Use Case Diagram}		}											
\end{flushleft}		
\blindtext
\begin{flushleft}
{\fontsize{14}{20}\selectfont\textbf{\*4.3 Dataflow Diagram}		}											
\end{flushleft}		
\blindtext \\

\setcounter{figure}{0}                   % 
\renewcommand{\thefigure}{4.\arabic{figure}}
\begin{figure}
    \centering
    \includegraphics[scale=0.5]{college_logo.JPG}
   \caption{Name of the figure}
    \end{figure}



\begin{flushleft}
{\fontsize{12}{12}\selectfont\textbf{\*4.3.1 SUB SECTION 3}		}											
\end{flushleft}		
\blindtext

\begin{flushleft}
{\fontsize{14}{20}\selectfont\textbf{\*4.4 SECTION 4}		}											
\end{flushleft}		
\blindtext

\begin{flushleft}
{\fontsize{12}{12}\selectfont\textbf{\*4.4.1 SUB SECTION 1}		}											
\end{flushleft}		
\blindtext
\renewcommand{\thefigure}{4.\arabic{figure}}
\begin{figure}
    \centering
    \includegraphics[scale=0.5]{Picture1.png}
   \caption{Name of the figure}
    \end{figure}
    \blindtext
\begin{flushleft}
{\fontsize{12}{12}\selectfont\textbf{\*4.4.2 SUB SECTION 2}		}											
\end{flushleft}		
\blindtext

\newpage
\onehalfspacing	
\begin{center}
{\fontsize{18pt}{18pt}\selectfont \textbf{CONCLUSION}}
\end{center}
\begin{flushleft}
    %{\fontsize{16pt}{16pt}\selectfont \textbf{}}
\end{flushleft}
\blindtext\\[2mm]
\newpage
\begin{center}

    \bibliographystyle{IEEEtran}
\bibliography{Reference_Paper}
\end{center}
\end{document}
